\chapter[2012 --- Gurdon et al.]{2012: John B. Gurdon, Shinya Yamanaka}
\label{ch:2012_Gurdon}

\section*{Main Contribution}

\textit{Discovered that mature cells can be reprogrammed to become pluripotent stem cells, opening new possibilities for regenerative medicine and disease modeling.}

\section{Official Citation}

for the discovery that mature cells can be reprogrammed to become pluripotent

\section{Background and Historical Context}

The 2012 Nobel Prize in Physiology or Medicine recognized John B. Gurdon and Shinya Yamanaka ``for the discovery that mature cells can be reprogrammed to become pluripotent.'' Their discoveries, separated by nearly half a century, fundamentally overturned a central dogma in developmental biology---that cell differentiation is irreversible---and opened up entirely new possibilities for regenerative medicine, disease modeling, and drug discovery.

For much of the twentieth century, the prevailing view in developmental biology held that cell differentiation was a one-way process. Once a cell had committed to a particular developmental fate---becoming a muscle cell, a nerve cell, a blood cell, or any other specialized cell type---there was no turning back. The differentiated state was thought to be stable and irreversible, maintained by the progressive loss of developmental potential and the acquisition of cell-type-specific patterns of gene expression. This view, which was deeply entrenched in the biological sciences, was supported by numerous observations that differentiated cells maintained their identity through many rounds of cell division and were not observed to spontaneously convert from one cell type to another. The molecular basis for this irreversibility was thought to involve the progressive loss of gene regulatory regions and the permanent silencing of genes associated with alternative cell fates. The prevailing assumption was that during differentiation, cells progressively restricted their developmental potential through irreversible changes in gene expression, ultimately arriving at a terminal, fixed state.

This dogma was challenged by a remarkable experiment conducted by John B. Gurdon in 1962. Gurdon, a British biologist working at the University of Oxford, transferred the nucleus of a differentiated intestinal epithelial cell from a tadpole of the African clawed frog \textit{Xenopus laevis} into an enucleated egg cell---an egg from which the nucleus had been removed. He then stimulated the reconstructed egg to divide and develop. The result was notable: some of the reconstructed eggs developed into normal, swimming tadpoles, demonstrating that the nucleus of a differentiated cell still contained all the genetic information necessary to direct the development of an entire organism.

This experiment, known as somatic cell nuclear transfer (SCNT), or ``nuclear cloning,'' provided the first direct evidence that cell differentiation did not involve the irreversible loss of genetic information. Rather, the differentiated state was maintained by epigenetic mechanisms---modifications to the DNA and its associated proteins that regulated gene expression without altering the underlying DNA sequence. The nucleus of a differentiated cell, when placed in the cytoplasm of an egg cell, could be reprogrammed by factors in the egg to revert to a totipotent state---a state from which all cell types in the body could be derived.

Gurdon's experiments were met with significant skepticism and criticism. The idea that a differentiated cell could be reprogrammed to direct the development of an entire organism challenged deeply held beliefs about the irreversibility of cell differentiation. Moreover, the efficiency of nuclear transfer was very low---only a small fraction of reconstructed eggs developed into normal tadpoles---and some critics questioned whether the successful results might be due to contamination with undifferentiated cells. However, Gurdon's findings were subsequently confirmed and extended by numerous laboratories, and his work laid the conceptual foundation for the field of cellular reprogramming.

\section{The Discovery/Contribution}

While Gurdon demonstrated that differentiated cells could be reprogrammed using factors present in the egg cytoplasm, the identity of these reprogramming factors remained unknown for nearly four decades. The search for the molecular mechanisms of reprogramming became one of a major questions in cell biology, with significant implications for the development of stem cell-based therapies.

The breakthrough came in 2006 from the laboratory of Shinya Yamanaka, a Japanese stem cell researcher at Kyoto University. Yamanaka had been studying the transcription factors that maintain the pluripotent state of embryonic stem cells (ESCs)---the stem cells derived from the inner cell mass of early embryos that have the potential to differentiate into virtually any cell type in the body. He hypothesized that a small number of key transcription factors might be sufficient to reprogram differentiated cells back to a pluripotent state.

Yamanaka and his colleague Kazutoshi Takahashi tested this hypothesis by systematically screening a panel of 24 genes that were known to be highly expressed in embryonic stem cells. They introduced each gene, or combinations of genes, into cultured mouse fibroblasts (connective tissue cells) using retroviral vectors, and then looked for the activation of a reporter gene that was specifically expressed in undifferentiated, pluripotent cells.

After extensive screening, Yamanaka and Takahashi identified just four transcription factors that were sufficient to reprogram mouse fibroblasts into pluripotent cells when expressed together. These four factors---Oct3/4, Sox2, Klf4, and c-Myc---were named ``Yamanaka factors,'' and the reprogrammed cells were called ``induced pluripotent stem cells'' (iPSCs). The iPSCs exhibited the key properties of embryonic stem cells: they expressed pluripotency markers, had the morphology of ESCs, could be differentiated into all three germ layers (ectoderm, mesoderm, and endoderm), and could contribute to chimeric mice when injected into blastocysts, including contributions to the germline.

The publication of Yamanaka's findings in \textit{Cell} in November 2006 was one of the most exciting moments in the history of biology. The ability to reprogram differentiated cells into a pluripotent state using just four transcription factors was a significant advance that opened up notable possibilities for regenerative medicine and disease research.

In 2007, Yamanaka's group, independently of a parallel effort by James Thomson's laboratory at the University of Wisconsin, reported the generation of human iPSCs. This achievement demonstrated that the reprogramming technology could be applied to human cells, raising the prospect of generating patient-specific stem cells for therapeutic purposes. The generation of human iPSCs also had the advantage of avoiding the ethical controversies associated with the use of human embryonic stem cells, as iPSCs could be derived from adult somatic cells without destroying embryos.

Subsequent research by Yamanaka and many other investigators has refined and improved the reprogramming process. Alternative methods for delivering the Yamanaka factors---including non-integrating viral vectors, plasmid DNA, proteins, and small molecules---have been developed to address safety concerns associated with the use of retroviral vectors, which can cause insertional mutagenesis. The discovery that some of the Yamanaka factors are dispensable, and that combinations of just three factors (Oct3/4, Sox2, and Klf4) can reprogram cells under certain conditions, has simplified the reprogramming process. Direct reprogramming---the conversion of one differentiated cell type to another without going through a pluripotent intermediate---has also been achieved, further demonstrating the plasticity of cellular identity.

The mechanisms underlying reprogramming have been gradually elucidated. Yamanaka factors work by binding to specific DNA sequences in the regulatory regions of genes involved in pluripotency and activating their expression. They also recruit chromatin-modifying enzymes that remodel the epigenetic landscape of differentiated cells, erasing cell-type-specific epigenetic marks and establishing the chromatin configuration characteristic of pluripotent cells. The reprogramming process is slow, taking approximately two to three weeks, and is highly inefficient---only a small fraction of cells exposed to the Yamanaka factors successfully reprogram to iPSCs. This inefficiency reflects the multiple epigenetic barriers that must be overcome during reprogramming.

\section{Impact on Modern Medicine}

The impact of Gurdon and Yamanaka's discoveries on modern medicine has been transformative. The development of iPSC technology has opened up new avenues for disease modeling, drug discovery, and cell-based therapies, and has the potential to revolutionize the practice of medicine.

One of a major applications of iPSC technology is in disease modeling. By generating iPSCs from patients with specific genetic diseases and differentiating them into the relevant cell types, researchers can create ``disease in a dish'' models that recapitulate key features of the human disease. These models can be used to study disease mechanisms, screen for therapeutic drugs, and develop personalized treatment strategies. iPSC-based disease models have been developed for a wide range of conditions, including neurodegenerative diseases (Alzheimer's disease, Parkinson's disease, amyotrophic lateral sclerosis), cardiovascular diseases, metabolic diseases, and rare genetic disorders.

The use of iPSCs for drug discovery and toxicology testing is also a major area of application. iPSC-derived cells can be used to test the efficacy and safety of drug candidates in human-relevant cell types, potentially reducing the need for animal testing and improving the predictability of clinical trial outcomes. The ability to generate large quantities of human cells from iPSCs has also enabled the development of high-throughput screening platforms for drug discovery.

In regenerative medicine, iPSCs hold great promise for the development of cell-based therapies for a wide range of diseases. Clinical trials involving iPSC-derived cells are underway for conditions including macular degeneration, Parkinson's disease, heart failure, and spinal cord injury. In 2014, a Japanese team led by Masayo Takahashi reported the first clinical application of iPSC-derived cells, transplanting iPSC-derived retinal pigment epithelial cells into a patient with age-related macular degeneration.

The generation of iPSCs from individual patients has also enabled the development of autologous cell therapies---therapies using the patient's own cells, which avoids the risk of immune rejection. This approach is particularly promising for conditions in which the patient's own cells are damaged or dysfunctional, such as macular degeneration, Parkinson's disease, and heart failure.

The ethical implications of iPSC technology are also significant. While iPSCs avoid some of the ethical concerns associated with embryonic stem cells, they raise new ethical questions about the potential for creating human embryos from iPSCs, the use of iPSCs for genetic modification of human embryos, and the equitable access to iPSC-based therapies. These ethical issues are the subject of ongoing discussion and regulation.

John B. Gurdon and Shinya Yamanaka were awarded the Nobel Prize in Physiology or Medicine in 2012 ``for the discovery that mature cells can be reprogrammed to become pluripotent.'' Their discoveries have fundamentally changed our understanding of cell differentiation and have opened up new possibilities for regenerative medicine, disease modeling, and drug discovery.

\section{Legacy}

The legacy of Gurdon and Yamanaka's work extends far beyond its immediate applications in research and medicine. Their discoveries have fundamentally changed our understanding of cellular identity and developmental potential, demonstrating that cell differentiation is not irreversible but can be reversed by the introduction of a small number of key transcription factors.

Gurdon's pioneering nuclear transfer experiments in the 1960s challenged the prevailing dogma of irreversible cell differentiation and provided the conceptual foundation for the field of cellular reprogramming. His work demonstrated that the genome of a differentiated cell retains the full complement of genetic information necessary to direct the development of an entire organism, and that this genetic information can be reactivated by factors in the egg cytoplasm.

Yamanaka's discovery of iPSCs provided a practical and powerful tool for cellular reprogramming that has been widely adopted by the scientific community. The development of iPSC technology has accelerated the pace of discovery in stem cell biology, regenerative medicine, and disease research, and has enabled the generation of patient-specific stem cells for therapeutic purposes.

Together, Gurdon and Yamanaka's work has created a new paradigm in biology---one in which cellular identity is viewed as plastic and malleable, rather than fixed and irreversible. This paradigm has inspired new approaches to the treatment of degenerative diseases, the study of human development, and the understanding of aging and age-related diseases. The legacy of their discoveries continues to shape the future of medicine and biology, and their contributions will remain central to these fields for generations to come.


\section{Modern Developments}

Gurdon and Yamanaka's reprogramming work has led to modern regenerative medicine. Induced pluripotent stem cells (iPSCs) are being used to model diseases, test drugs, and develop cell therapies. Understanding of reprogramming factors has led to direct conversion of fibroblasts into neurons, cardiomyocytes, and other cell types. Clinical trials using iPSC-derived retinal pigment epithelium for macular degeneration are underway. Understanding of epigenetic reprogramming has informed development of cancer therapies.
